
\section*{Глава 4. Оформление итогов работы}
\addcontentsline{toc}{section}{Глава 4. Оформление итогов работы}
\stepcounter{section}

\subsection{Создание Git-репозитория с кодом проекта}

Для отслеживания истории изменений и возможности совместной разработки была использована распределенная система контроля версий \textbf{Git}, а в качестве 
платформы для удаленного хостинга репозитория был выбран сервис \textbf{GitHub}. Архитектура хранения кода реализована по паттерну 
монорепозитория, где исходный код разделен на две логические директории: серверную (\texttt{backend}) и 
клиентскую (\texttt{frontend}) части. Данный подход упрощает синхронизацию версий API и клиентского интерфейса.

Особое внимание было уделено информационной безопасности. Был сформирован файл исключений \texttt{.gitignore}, предотвращающий случайную утечку конфиденциальных 
данных в публичный доступ. Из индексации Git были исключены локальные конфигурационные файлы по типу .env, объемные директории зависимостей и каталоги с загружаемыми 
пользовательскими медиафайлами. Для удобства навигации и проверки проекта в корень репозитория был добавлен файл \texttt{README.md}, содержащий техническое описание 
стека, инструкции по локальному запуску и ссылки на развернутое приложение.

\subsection{Деплой приложения на хостинг}

Для деплоя продукта была спроектирована облачная инфраструктура, состоящая из базы данных, серверного API и клиентского веб-приложения.

В качестве СУБД был использован облачный кластер \textbf{PostgreSQL}, развернутый на платформе \textbf{Render}. Бэкенд также был опубликован на платформе Render в формате веб-сервиса. 
Процесс развертывания сервера включал настройку CI/CD с привязкой к репозиторию на GitHub. Помимо этого производилась настройка переменных окружения для безопасного 
подключения к удаленной БД, а также настройка политик CORS для разрешения кросс-доменных HTTP-запросов от клиента.

Развертывание клиентской части было осуществлено на платформе \textbf{Vercel}, оптимизированной для работы с современными JavaScript-фреймворками. Сборка React-приложения 
выполнялась инструментом Vite на серверах платформы Vercel. На этапе сборки в приложение была успешно интегрирована переменная, которая направляла
клиентские запросы на рабочий домен бэкенда.

В результате успешного деплоя продукт стал доступен конечным пользователям через глобальную сеть Интернет~\cite{deploy_host}.

\subsection{Оформление отчета о проделанной работе}

Были проведены систематизация результатов и формирование итогового отчета, поскольку документирование является неотъемлемым этапом разработки любого проекта.

В процессе оформления курсового проекта были выполнены пункты, описанные ниже.
\begin{enumerate}
    \item \textbf{Аналитическое обоснование:} структурирован обзор предметной области, проведен сравнительный анализ существующих решений на современном рынке и обоснован 
	выбор стека технологий.
    \item \textbf{Визуализация архитектуры:} разработаны графические схемы, включая UML-диаграмму вариантов использования для формализации ролевой модели и 
	детализированную инфологическую ER-диаграмму реляционной базы данных с описанием физических схем таблиц.
    \item \textbf{Прототипирование:} составлены структурные Wireframe-макеты ключевых пользовательских интерфейсов.
    \item \textbf{Текстовое оформление:} пояснительная записка приведена в соответствие с требованиями академических стандартов и ГОСТ с помощью набора \textbf{LaTeX}.
\end{enumerate}

Подготовленная отчетная документация полностью отражает процесс разработки веб-платформы маркетплейса от первоначального этапа сбора требований до вывода продукта в 
эксплуатацию, подтверждая успешное выполнение всех поставленных задач курсового проекта.

